%REVIEW-ADDR: W1-zvf-tautology
%REVIEW-ADDR: Q1-zvf-formalization
\section{Formal Definition of the Zero-Variance Fraction, Non-Tautology, and Scope}
\label{sec:appendix:zvf-formalization}

\subsection{Formal Definition}

Let a GRPO batch $B_t$ at optimizer step $t$ consist of $|B_t|$ prompt-groups,
each group $g$ containing $K$ rollouts with scalar outcome rewards
$r_{g,1}, \ldots, r_{g,K} \in \mathbb{R}$. We define the \emph{Zero-Variance
Fraction} (ZVF) as
\begin{equation}
\label{eq:zvf}
\mathrm{ZVF}_t \;=\; \frac{1}{|B_t|} \sum_{g \in B_t}
  \mathbb{1}\!\left[\;\widehat{\mathrm{Var}}(r_{g,1}, \ldots, r_{g,K}) \le \varepsilon\;\right],
\end{equation}
where $\widehat{\mathrm{Var}}$ is the unbiased sample variance and
$\varepsilon$ is a small numerical tolerance
(we use $\varepsilon = 10^{-6}$ for rewards normalized to $[0,1]$).
Intuitively, $\mathrm{ZVF}_t$ is the fraction of groups at step $t$ whose
rollouts all collapsed to the same outcome; those groups contribute zero
group-relative advantage and thus zero gradient signal to the GRPO update.

\subsection{Why ZVF Is Not a Tautological Re-expression of Mean Reward}

A reviewer may reasonably suspect that under binary outcome rewards
$r \in \{0,1\}$ and group-relative advantages, $\mathrm{ZVF}_t$ is a direct
function of the batch mean reward $\bar{r}_t$: if every group succeeds
(or fails) uniformly, $\mathrm{ZVF}_t = 1$. This is true at the extremes
$\bar{r}_t \in \{0,1\}$ but does \emph{not} hold at intermediate values.
For $K$ i.i.d.~Bernoulli rollouts with prompt-specific success probability
$p_g$, the expected ZVF under a null independence model is
\begin{equation}
\mathbb{E}[\mathrm{ZVF}_t] \;=\; \mathbb{E}_{p_g}\!\left[\,p_g^{K} + (1-p_g)^{K}\,\right].
\end{equation}
The right-hand side depends on the full prompt-difficulty distribution
$\{p_g\}$, not only on its first moment. Two populations can therefore
share the same mean reward $\bar r = 0.5$ yet have very different ZVF:
a bimodal population with $p_g \in \{0,1\}$ yields $\mathrm{ZVF}=1$,
while a uniform-hard population with all $p_g = 0.5$ and $K=8$ yields
$\mathrm{ZVF} \approx 2 \cdot (0.5)^8 \approx 0.008$. The substantive
claim is thus not ``high mean reward implies high ZVF,'' but rather that
ZVF is sensitive to \emph{how} success mass is distributed across prompts.

\subsection{What the Released Artifact Establishes}

The main Qwen3-8B/frontier GSM8K corpus bundled with this paper contains
aggregate reward trajectories, held-out checkpoint summaries, and the
cross-framework parser specification in Appendix~\ref{sec:appendix:zvf-pipeline},
but \emph{not} the per-group step-level rollout logs for those frontier runs.
For that corpus we therefore make no causal partial-correlation claim. To test
the partial-correlation question directly we instead ran a fully instrumented,
ungated small-scale validation (next subsection), which \emph{does} log
per-step ZVF together with batch-mean reward, policy entropy, and advantage
variance, and we report its measured partial correlation there.

The narrower empirical claim made in the revised paper is the one directly
supported by the released evidence: in outcome-reward GRPO, high ZVF
coincides with the disappearance of within-group contrast, and those are
exactly the runs where reward traces plateau or regress because the
group-relative advantage has collapsed to zero. This is the operational
meaning of ZVF in the present benchmark.

\subsection{Measured Validation of the ZVF Formalism (Small-Scale and Frontier)}
\label{sec:appendix:zvf-measured}

To test the formalism end-to-end with fully logged rollouts we ran an
instrumented GRPO sweep on the ungated \texttt{Qwen2.5-0.5B} with a verifiable
binary correctness reward on two-operand addition (so within-group reward
variance is real), at group sizes $G \in \{2,4,8,16\}$ across three seeds, $40$
optimizer steps each ($12$ runs, $480$ logged steps; driver
\texttt{experiments/modal/modal\_groupsize\_zvf\_sweep.py}, artifacts
\texttt{experiments/results/groupsize\_zvf\_sweep.\{json,tsv\}} and
\texttt{zvf\_partial\_correlations.tsv}). This is a deliberately small,
ungated probe of the \emph{mechanism}, not a restatement of the
Qwen3-8B/GSM8K results.

\paragraph{The closed-form $\mathrm{ZVF}(p,G)=p^{G}+(1-p)^{G}$ tracks the data.}
Across the $480$ logged steps, the per-step observed ZVF correlates
$r=0.92$ with the Bernoulli-null prediction $p^{G}+(1-p)^{G}$ evaluated at the
step's batch-mean reward $p$, and mean ZVF decreases monotonically in $G$
($0.84, 0.76, 0.69, 0.63$ at $G=2,4,8,16$), as the closed form requires. The
\emph{absolute} ZVF exceeds the i.i.d.\ prediction (e.g.\ $0.69$ vs.\ a predicted
$0.33$ at $G{=}8$), consistent with the low-temperature/correlated-completion
boundary condition noted above: the $G$ rollouts are not independent draws, so
the independence null is a lower bound on observed ZVF.

\paragraph{The closed form captures the qualitative structure at frontier scale,
with a circularity caveat.} We measured ZVF directly on the frontier model by
sampling \texttt{Qwen3-8B} on a held-out GSM8K slice via the Tinker API
($G{=}8$, temperature~$1.0$, $200$ problems $\times$ three seeds; driver
\texttt{experiments/tinker\_direct\_eval.py}, artifacts
\texttt{experiments/results/tinker\_gsm8k\_zvf\_*.json}). Mean ZVF is $0.16$,
concentrated on the saturated tails of the difficulty distribution ($12.7\%$
all-correct, $3.2\%$ all-wrong, $84\%$ mixed)---the qualitative structure
$\mathrm{ZVF}(p,G)$ predicts. We caution that a per-problem Pearson correlation
between observed ZVF and $p^{G}+(1-p)^{G}$ is \emph{mechanically} inflated: the
per-problem observed ZVF is just the binary all-agree indicator
$\mathbf{1}[p \in \{0,1\}]$ and the closed form is a deterministic function of
the same $p$, so the large $r \approx 0.93$ chiefly reflects their shared
dependence on $p$'s extremity (the rank correlation is weaker,
$\rho \approx 0.65$). The i.i.d.-Bernoulli null is moreover only a coarse level
estimate---here it \emph{over}-predicts mean ZVF ($\approx 0.28$ vs.\ observed
$0.16$), the opposite-sign residual to the low-temperature $0.5$B probe. We thus
read this as a consistency check on the \emph{shape} of the diagnostic across
scales, not a precise quantitative fit.

\paragraph{ZVF vs.\ mean reward (W1): what we can and cannot claim.} On the
logged per-step rows ZVF is strongly correlated with batch-mean reward
($r=+0.74$): more all-correct groups means more zero-variance groups, as
expected. We deliberately do \emph{not} report a partial correlation
``controlling for advantage variance'': under binary outcome rewards the
within-group advantage variance is a deterministic transform of ZVF (it is zero
\emph{iff} the group is zero-variance), so partialling it out regresses ZVF on a
function of itself---a circular control. Nor do we attach a step-level
$p$-value: the $480$ rows are $12$ runs $\times 40$ \emph{autocorrelated}
optimizer steps (lag-1 $\approx 0.9$), not independent observations. The
defensible statement is the descriptive one in the main text: high ZVF coincides
with the disappearance of within-group contrast, and across runs in this
near-solved regime ZVF does not track final accuracy ($r\approx 0.09$,
$n{=}12$)---the expected behaviour once the task is essentially solved.

\subsection{Boundary Conditions and Scope}

We state explicitly when ZVF is \emph{not} informative, partly in response
to the tautology concern:
\begin{itemize}[leftmargin=1.2em,itemsep=0pt]
  \item \textbf{Dense or continuous rewards.} When rewards are continuous
        (e.g.,~process-reward models or graded code-execution partial credit),
        $\widehat{\mathrm{Var}}$ is almost surely nonzero and
        $\mathrm{ZVF} \to 0$. ZVF degenerates and must be replaced by a
        per-step-reward-variance diagnostic or the surrogates discussed in
        Section~\ref{sec:extended-related-work}.
  \item \textbf{$K=1$.} Group variance is undefined and ZVF is trivially $1$.
  \item \textbf{SFT-saturated baselines.} When the policy already solves
        the task ($\bar p_g \to 1$ for essentially every prompt),
        $\mathrm{ZVF} \to 1$ and the diagnostic loses discriminative power;
        in this regime RL no longer provides useful signal anyway.
  \item \textbf{Non-outcome-reward RL.} For DPO- or regression-style
        training there are no rollout groups and ZVF is ill-defined.
\end{itemize}
Within its intended scope -- outcome-reward GRPO on verifiable tasks with
$K \ge 4$ and non-saturated reward support -- ZVF is a cheap, model-agnostic
early-warning diagnostic. Outside that scope, the paper now recommends
explicit surrogates instead of over-claiming portability.
